\section{Instalación y configuración de PostgreSQL}

\subsection{Descarga de la imagen}

PostgreSQL se desplegó dentro de un contenedor, por lo que no fue necesario
instalar el servidor directamente sobre Fedora. Se eligió la etiqueta
\texttt{postgres:18} en lugar de \texttt{latest}; de esta manera se mantiene
fija la versión principal y el procedimiento puede repetirse con resultados
previsibles.

\begin{lstlisting}[style=terminal]
$ docker pull postgres:18
\end{lstlisting}

La descarga produjo el resumen SHA-256 de la imagen y registró PostgreSQL 18.6
como versión disponible en el momento de la práctica.

\subsection{Credencial y creación del contenedor}

La contraseña se generó aleatoriamente en un archivo temporal. El uso de
\texttt{--env-file} impidió que el valor quedara visible en el comando o en la
captura de pantalla.

\begin{lstlisting}[style=terminal]
$ umask 077
$ printf 'POSTGRES_PASSWORD=' > /tmp/nipgg_practica01_postgres.env
$ openssl rand -base64 24 | tr -d '\n' >> /tmp/nipgg_practica01_postgres.env
$ printf '\n' >> /tmp/nipgg_practica01_postgres.env
\end{lstlisting}

Se creó el contenedor con un nombre descriptivo, un volumen persistente y un
mapeo de red restringido a la interfaz local. La restricción
\texttt{127.0.0.1} evita exponer el servidor directamente a otros equipos de la
red.

\begin{lstlisting}[style=terminal]
$ docker run -d \
    --name postgres-practica01 \
    --env-file /tmp/nipgg_practica01_postgres.env \
    -p 127.0.0.1:5432:5432 \
    -v postgres_practica01_data:/var/lib/postgresql \
    postgres:18
\end{lstlisting}

\evidencia{03_evidencia.png}{Imagen, contenedor, puerto local y volumen persistente de PostgreSQL.}

\subsection{Comprobación con psql}

Los registros del contenedor indicaron que el sistema estaba listo para
aceptar conexiones. A continuación se abrió \texttt{psql} dentro del propio
contenedor y se consultaron la base activa, el usuario y la versión del
servidor:

\begin{lstlisting}[style=terminal]
$ docker logs --tail 12 postgres-practica01
$ docker exec postgres-practica01 psql -U postgres -d postgres \
    -c 'SELECT current_database(), current_user, version();'
\end{lstlisting}

\evidencia{04_evidencia.png}{Consulta de validación ejecutada con \texttt{psql} sobre PostgreSQL 18.6.}
